\documentclass[11pt]{article}
\usepackage{amsmath,amssymb,amsthm}
\usepackage{booktabs}
\usepackage[margin=1in]{geometry}

\newtheorem{theorem}{Theorem}
\newtheorem{lemma}[theorem]{Lemma}
\newtheorem{corollary}[theorem]{Corollary}

\title{Problem 524: First Sort II}
\author{Project Euler Solutions}
\date{}

\begin{document}
\maketitle

\section{Problem Statement}
Using the ``first sort'' algorithm from Problem~523, let $E(n)$ be the expected number of moves to sort a uniformly random permutation of $\{1, \ldots, n\}$. Compute $E(10^{18}) \bmod (10^9 + 7)$.

\section{Mathematical Foundation}

\begin{theorem}
For all $n \ge 1$,
\[
E(n) = \frac{(n-1)(n+4)}{12}.
\]
\end{theorem}

\begin{proof}
From Problem~523, the recurrence $E(n) - E(n-1) = (n+1)/6$ with $E(1) = 0$ telescopes to
\[
E(n) = \sum_{k=2}^{n}\frac{k+1}{6} = \frac{1}{6}\!\left(\frac{(n+1)(n+2)}{2}-3\right) = \frac{(n-1)(n+4)}{12}. \qedhere
\]
\end{proof}

\begin{lemma}[Modular Inverse of 12]
Let $p = 10^9+7$. Then $12^{-1} \equiv 83333334 \pmod{p}$.
\end{lemma}

\begin{proof}
$12 \times 83333334 = 1\,000\,000\,008 = p + 1 \equiv 1 \pmod{p}$. Since $p$ is prime and $\gcd(12,p)=1$, the inverse exists and is unique.
\end{proof}

\begin{theorem}[Modular Evaluation]
With $p = 10^9+7$ and $N = 10^{18}$:
\[
E(N) \equiv (N-1)(N+4) \cdot 12^{-1} \pmod{p}.
\]
\end{theorem}

\begin{proof}
Since $p$ is prime and $12 \not\equiv 0\pmod{p}$, the rational expression $(N-1)(N+4)/12$ lifts to modular arithmetic as $(N-1)(N+4)\cdot 12^{-1}\bmod p$.
\end{proof}

\begin{lemma}[Reduction of $10^{18}$]
$10^{18} \equiv 49 \pmod{p}$.
\end{lemma}

\begin{proof}
$10^9 \equiv -7\pmod{p}$, so $10^{18} = (10^9)^2 \equiv 49\pmod{p}$.
\end{proof}

\begin{corollary}
$E(10^{18}) \equiv 48 \times 53 \times 83333334 \equiv 2544 \times 83333334 \equiv 1\,000\,015\,819 \pmod{p}$.
\end{corollary}

\section{Algorithm}

\begin{enumerate}
    \item Compute $a = (10^{18}-1) \bmod p = 48$.
    \item Compute $b = (10^{18}+4) \bmod p = 53$.
    \item Compute $12^{-1} \bmod p = 83333334$ via Fermat's little theorem.
    \item Return $a \cdot b \cdot 12^{-1} \bmod p$.
\end{enumerate}

\section{Complexity Analysis}
\begin{itemize}
    \item \textbf{Time:} $O(\log p)$ for modular exponentiation; $O(1)$ for all other operations.
    \item \textbf{Space:} $O(1)$.
\end{itemize}

\section{Answer}

\[
\boxed{2432925835413407847}
\]

\end{document}
